% Target: rmp-workbench-iii-peps-renormalization-one
% Formula: Four PEPS tensors occupy the edges of one renormalization plaquette.
% Ink: Kernel flat frame; eight virtual half-indices cross at the corners
%   between four edge-centred glyphless sites with angled physical legs.
% Boundary: Eight external half-edges and four physical legs remain open.
% Source: author source ImagesReview Section III/ImagesReview Section III.tex
%   lines 622-630; archival output PEPSrenorm1.pdf
% Capabilities: kernel, four-site-plaquette, renormalization
\[
  \tenkzkernel
  \begin{tenkz}[rows={wire}, cols=1, bonds=none]
    % Each edge site has two open virtual indices. Ordinary index lines
    % cross without inventing an over/under convention at the corners.
    \tnwire{open w}{PN}
    \tnwire{PN}{open e}
    \tnwire{open w}{PS}
    \tnwire{PS}{open e}
    \tnwire{open n}{PW}
    \tnwire{PW}{open s}
    \tnwire{open n}{PE}
    \tnwire{PE}{open s}
    % the four edge-centred PEPS sites are glyphless; each shows only its
    % 50-degree physical leg, as the source draws
    \tn[skin=none, at=0.375 n of (1,1), name=PN, ports={50:physical}]{}
    \tn[skin=none, at=0.375 s of (1,1), name=PS, ports={50:physical}]{}
    \tn[skin=none, at=0.375 w of (1,1), name=PW, ports={50:physical}]{}
    \tn[skin=none, at=0.375 e of (1,1), name=PE, ports={50:physical}]{}
  \end{tenkz}
\]
